\section{Considerações Finais}
\label{sec:conclusao}

A associação entre preços agrícolas e conversão da vegetação em Mato Grosso depende do universo observado e da forma de representar suas diferenças produtivas e temporais. Ao incorporar o histórico necessário às perdas, os gradientes positivos de preço×capacidade diminuem antes mesmo da inclusão de H1b. No núcleo com perdas, os intervalos de H1a abrangem zero para soja e milho. Os modificadores preço×perda também são imprecisos, com sinal positivo na soja e próximo de zero no milho. Não há, portanto, suporte suficientemente estável para confirmar as duas proposições de H1 de modo geral, sem que isso demonstre ausência das relações investigadas.

As comparações complementares mostram por que a média do painel não esgota a questão. A diferença do gradiente preço×capacidade entre agricultura e pastagem é negativa no milho, com precisão sensível à inferência. Os regimes e o contraste do milho no território alcançado pela Moratória apresentam rejeições sob HAC que não se reproduzem de forma uniforme nos agrupamentos municipais. Esses resultados preservam a possibilidade de heterogeneidade, mas a concentração de informação e as janelas curtas impedem convertê-los em uma resposta geral ou em efeitos das políticas.

A contribuição do artigo está em distinguir potencial produtivo, desempenho realizado e conversão por destino numa comparação municipal comum, mostrando como composição, calendário e precisão condicionam sua leitura. A participação da base natural no processo econômico, a importância das condições locais e a dimensão temporal das decisões orientam essa construção. O desenho não observa expectativas individuais nem demonstra que perdas tenham restringido recursos ou motivado expansão compensatória. Da mesma forma, o alcance potencial da Moratória e a cobertura florestal inicial não substituem cumprimento comercial ou obrigação jurídica por imóvel.

O avanço dessa interpretação depende de informações que aproximem perdas, decisões e restrições no mesmo território: séries por safra, limites históricos compatíveis e registros de imóveis, financiamento e cumprimento institucional. No recorte atual, a evidência recomenda distinguir as alternativas de uso da terra e preservar a incerteza dos contrastes, em vez de deduzir maior abertura do crescimento da oferta ou eficácia institucional da simples mudança de coeficientes entre períodos.
